% Options for packages loaded elsewhere
% Options for packages loaded elsewhere
\PassOptionsToPackage{unicode}{hyperref}
\PassOptionsToPackage{hyphens}{url}
\PassOptionsToPackage{dvipsnames,svgnames,x11names}{xcolor}
%
\documentclass[
  spanish,
  10pt,
  a4paper,
  DIV=11,
  numbers=noendperiod]{scrartcl}
\usepackage{xcolor}
\usepackage[top=18mm,bottom=18mm,left=18mm,right=18mm]{geometry}
\usepackage{amsmath,amssymb}
\setcounter{secnumdepth}{-\maxdimen} % remove section numbering
\usepackage{iftex}
\ifPDFTeX
  \usepackage[T1]{fontenc}
  \usepackage[utf8]{inputenc}
  \usepackage{textcomp} % provide euro and other symbols
\else % if luatex or xetex
  \usepackage{unicode-math} % this also loads fontspec
  \defaultfontfeatures{Scale=MatchLowercase}
  \defaultfontfeatures[\rmfamily]{Ligatures=TeX,Scale=1}
\fi
\usepackage{lmodern}
\ifPDFTeX\else
  % xetex/luatex font selection
\fi
% Use upquote if available, for straight quotes in verbatim environments
\IfFileExists{upquote.sty}{\usepackage{upquote}}{}
\IfFileExists{microtype.sty}{% use microtype if available
  \usepackage[]{microtype}
  \UseMicrotypeSet[protrusion]{basicmath} % disable protrusion for tt fonts
}{}
\makeatletter
\@ifundefined{KOMAClassName}{% if non-KOMA class
  \IfFileExists{parskip.sty}{%
    \usepackage{parskip}
  }{% else
    \setlength{\parindent}{0pt}
    \setlength{\parskip}{6pt plus 2pt minus 1pt}}
}{% if KOMA class
  \KOMAoptions{parskip=half}}
\makeatother
% Make \paragraph and \subparagraph free-standing
\makeatletter
\ifx\paragraph\undefined\else
  \let\oldparagraph\paragraph
  \renewcommand{\paragraph}{
    \@ifstar
      \xxxParagraphStar
      \xxxParagraphNoStar
  }
  \newcommand{\xxxParagraphStar}[1]{\oldparagraph*{#1}\mbox{}}
  \newcommand{\xxxParagraphNoStar}[1]{\oldparagraph{#1}\mbox{}}
\fi
\ifx\subparagraph\undefined\else
  \let\oldsubparagraph\subparagraph
  \renewcommand{\subparagraph}{
    \@ifstar
      \xxxSubParagraphStar
      \xxxSubParagraphNoStar
  }
  \newcommand{\xxxSubParagraphStar}[1]{\oldsubparagraph*{#1}\mbox{}}
  \newcommand{\xxxSubParagraphNoStar}[1]{\oldsubparagraph{#1}\mbox{}}
\fi
\makeatother


\usepackage{longtable,booktabs,array}
\newcounter{none} % for unnumbered tables
\usepackage{calc} % for calculating minipage widths
% Correct order of tables after \paragraph or \subparagraph
\usepackage{etoolbox}
\makeatletter
\patchcmd\longtable{\par}{\if@noskipsec\mbox{}\fi\par}{}{}
\makeatother
% Allow footnotes in longtable head/foot
\IfFileExists{footnotehyper.sty}{\usepackage{footnotehyper}}{\usepackage{footnote}}
\makesavenoteenv{longtable}
\usepackage{graphicx}
\makeatletter
\newsavebox\pandoc@box
\newcommand*\pandocbounded[1]{% scales image to fit in text height/width
  \sbox\pandoc@box{#1}%
  \Gscale@div\@tempa{\textheight}{\dimexpr\ht\pandoc@box+\dp\pandoc@box\relax}%
  \Gscale@div\@tempb{\linewidth}{\wd\pandoc@box}%
  \ifdim\@tempb\p@<\@tempa\p@\let\@tempa\@tempb\fi% select the smaller of both
  \ifdim\@tempa\p@<\p@\scalebox{\@tempa}{\usebox\pandoc@box}%
  \else\usebox{\pandoc@box}%
  \fi%
}
% Set default figure placement to htbp
\def\fps@figure{htbp}
\makeatother



\ifLuaTeX
\usepackage[bidi=basic,shorthands=off]{babel}
\else
\usepackage[bidi=default,shorthands=off]{babel}
\fi
\ifLuaTeX
  \usepackage{selnolig} % disable illegal ligatures
\fi


\setlength{\emergencystretch}{3em} % prevent overfull lines

\providecommand{\tightlist}{%
  \setlength{\itemsep}{0pt}\setlength{\parskip}{0pt}}



 


\KOMAoption{captions}{tableheading}
\makeatletter
\@ifpackageloaded{caption}{}{\usepackage{caption}}
\AtBeginDocument{%
\ifdefined\contentsname
  \renewcommand*\contentsname{Tabla de contenidos}
\else
  \newcommand\contentsname{Tabla de contenidos}
\fi
\ifdefined\listfigurename
  \renewcommand*\listfigurename{Listado de Figuras}
\else
  \newcommand\listfigurename{Listado de Figuras}
\fi
\ifdefined\listtablename
  \renewcommand*\listtablename{Listado de Tablas}
\else
  \newcommand\listtablename{Listado de Tablas}
\fi
\ifdefined\figurename
  \renewcommand*\figurename{Figura}
\else
  \newcommand\figurename{Figura}
\fi
\ifdefined\tablename
  \renewcommand*\tablename{Tabla}
\else
  \newcommand\tablename{Tabla}
\fi
}
\@ifpackageloaded{float}{}{\usepackage{float}}
\floatstyle{ruled}
\@ifundefined{c@chapter}{\newfloat{codelisting}{h}{lop}}{\newfloat{codelisting}{h}{lop}[chapter]}
\floatname{codelisting}{Listado}
\newcommand*\listoflistings{\listof{codelisting}{Listado de Listados}}
\makeatother
\makeatletter
\makeatother
\makeatletter
\@ifpackageloaded{caption}{}{\usepackage{caption}}
\@ifpackageloaded{subcaption}{}{\usepackage{subcaption}}
\makeatother
\usepackage{bookmark}
\IfFileExists{xurl.sty}{\usepackage{xurl}}{} % add URL line breaks if available
\urlstyle{same}
\hypersetup{
  pdftitle={Informe histórico de mesa electoral},
  pdfauthor={Análisis electoral de Fuencarral-El Pardo},
  pdflang={es},
  colorlinks=true,
  linkcolor={blue},
  filecolor={Maroon},
  citecolor={Blue},
  urlcolor={Blue},
  pdfcreator={LaTeX via pandoc}}


\title{Informe histórico de mesa electoral}
\usepackage{etoolbox}
\makeatletter
\providecommand{\subtitle}[1]{% add subtitle to \maketitle
  \apptocmd{\@title}{\par {\large #1 \par}}{}{}
}
\makeatother
\subtitle{Elecciones generales al Congreso de los Diputados}
\author{Análisis electoral de Fuencarral-El Pardo}
\date{30 de octubre de 2026}
\begin{document}
\maketitle

\renewcommand*\contentsname{Tabla de contenidos}
{
\setcounter{tocdepth}{2}
\tableofcontents
}

\textbf{Mesa analizada:} 28-079-08-056-A \n**Ámbito:** Fuencarral-El
Pardo \n**Periodo disponible:** 03/1996 a 07/2023 \n**Número de
convocatorias:** 10

\section{Lectura rápida}\label{lectura-ruxe1pida}

En la convocatoria de \textbf{07/2023}, la mesa registró \textbf{504
electores}, una participación del \textbf{78,2 \%} y una abstención del
\textbf{21,8 \%}. La fuerza más votada fue \textbf{PP}, con \textbf{217
votos}. La participación disminuyó 5,6 puntos respecto de la
convocatoria anterior.

\section{Evolución del censo y la
participación}\label{evoluciuxf3n-del-censo-y-la-participaciuxf3n}

La comparación histórica debe interpretarse con cautela. Aunque el
código de mesa sea el mismo, las secciones electorales pueden dividirse,
fusionarse o modificar sus límites. Por ello, un cambio brusco de censo
puede indicar una alteración de la composición territorial de la mesa y
no únicamente un cambio demográfico.

\begin{figure}[H]

{\centering \includegraphics[width=1\linewidth,height=\textheight,keepaspectratio]{informe_mesa_files/figure-pdf/grafico-censo-1.pdf}

}

\caption{Evolución del censo electoral de la mesa.}

\end{figure}%

\begin{figure}[H]

{\centering \includegraphics[width=1\linewidth,height=\textheight,keepaspectratio]{informe_mesa_files/figure-pdf/grafico-participacion-1.pdf}

}

\caption{Participación y abstención sobre el censo electoral.}

\end{figure}%

\section{Evolución del voto}\label{evoluciuxf3n-del-voto}

\begin{figure}[H]

{\centering \includegraphics[width=1\linewidth,height=\textheight,keepaspectratio]{informe_mesa_files/figure-pdf/grafico-votos-1.pdf}

}

\caption{Número de votos por familia política y convocatoria.}

\end{figure}%

\begin{figure}[H]

{\centering \includegraphics[width=1\linewidth,height=\textheight,keepaspectratio]{informe_mesa_files/figure-pdf/grafico-porcentajes-1.pdf}

}

\caption{Porcentaje de cada familia sobre el total de votos a
candidaturas.}

\end{figure}%

\section{Última elección}\label{uxfaltima-elecciuxf3n}

\begin{figure}[H]

{\centering \includegraphics[width=1\linewidth,height=\textheight,keepaspectratio]{informe_mesa_files/figure-pdf/semitarta-1.pdf}

}

\caption{Distribución del voto a candidaturas en la última elección
disponible.}

\end{figure}%

\begin{longtable}[]{@{}lrr@{}}
\caption{Resultados por familias políticas en 07/2023}\tabularnewline
\toprule\noalign{}
Fuerza & Votos & Porcentaje \\
\midrule\noalign{}
\endfirsthead
\toprule\noalign{}
Fuerza & Votos & Porcentaje \\
\midrule\noalign{}
\endhead
\bottomrule\noalign{}
\endlastfoot
PP & 217 & 55,2\% \\
PSOE & 81 & 20,6\% \\
VOX & 49 & 12,5\% \\
IU / Podemos / Más País / Sumar & 45 & 11,5\% \\
Otros & 1 & 0,3\% \\
UPyD / Ciudadanos & 0 & 0,0\% \\
\end{longtable}

\section{Notas metodológicas}\label{notas-metodoluxf3gicas}

\begin{itemize}
\tightlist
\item
  Los porcentajes de las fuerzas políticas se calculan sobre los
  \textbf{votos a candidaturas}, sin incluir votos en blanco ni votos
  nulos.
\item
  \texttt{IU\_PODEMOS\_MAS\_PAIS\_SUMAR} agrupa las candidaturas
  identificadas como IU, Podemos, Más País, Más Madrid y Sumar.
\item
  \texttt{UPYD\_CIUDADANOS} agrupa las candidaturas identificadas como
  UPyD y Ciudadanos.
\item
  \texttt{OTROS} incluye el resto de candidaturas, entre ellas PACMA.
\item
  La continuidad del identificador no garantiza que la delimitación
  territorial y la composición del electorado hayan permanecido
  constantes.
\end{itemize}




\end{document}
